% Two distinct value interfaces for conditional action generation.
\begin{tikzpicture}[
  x=\linewidth,y=1mm,
  font=\figurefont\fontsize{8}{10}\selectfont,
  box/.style={draw=BookRule,fill=BookPaper,line width=.8pt,
    text width=.23\linewidth,minimum height=10mm,align=center,inner sep=1mm},
  flow/.style={-{Latex[length=1.8mm]},draw=BookInk,line width=.9pt},
  gradient/.style={-{Latex[length=1.8mm]},draw=BookAmber,line width=.8pt,dashed}
]
  \node[anchor=west] at (.01,12) {(a) 价值驱动生成器训练};
  \node[box,fill=BookTeal!8] (generator) at (.14,0) {条件生成器};
  \node[box] (candidate) at (.5,0) {生成动作$a$};
  \node[box] (value) at (.86,0) {价值$Q(s,a)$};
  \draw[flow] (generator) -- (candidate);
  \draw[flow] (candidate) -- (value);
  \draw[gradient] (value.south) -- (.86,-14)
    -- node[below] {价值梯度经生成过程回传} (.14,-14) -- (generator.south);
  \node[anchor=west] at (.01,-30) {(b) 行为生成后的价值筛选};
  \node[box,fill=BookTeal!8] (behavior) at (.14,-43) {行为生成器};
  \node[box] (candidates) at (.5,-43) {$a_1,\ldots,a_m$};
  \node[box,fill=BookBlue!8] (selection) at (.86,-43) {加权重采样\\或取最大};
  \node[box] (critic) at (.5,-66) {价值评分};
  \node[box] (action) at (.86,-66) {最终动作};
  \draw[flow] (behavior) -- (candidates);
  \draw[flow] (candidates) -- (selection);
  \draw[flow] (candidates) -- (critic);
  \draw[flow] (critic.east) -- (.7,-66) -- (.7,-55) -- (selection.south west);
  \draw[flow] (selection) -- (action);
\end{tikzpicture}
